\section{Probe definitions and configuration}
\label{app:probes}

Table~\ref{tab:config} records the constants that were frozen before expert-set scoring.
Lengths are fractions of the bounding-box diagonal after normalisation.
The genus feature passed to the classifier is \(\log(1+g)\) with \(g\) as in Equation~\eqref{eq:genus}.
Every other non-negative feature is passed through \(\log(1+x)\) and then standardised on the training split.
Standardisation, class weights and the decision threshold are part of the frozen predictor, not of the mesh probes.
The probes themselves have no fitted parameters.
Sampling uses seed 0.

\begin{table}[t]
  \centering
  \caption{Frozen probe and classifier configuration.}
  \label{tab:config}
  \small
  \begin{tabular}{lp{8.6cm}}
    \toprule
    Item & Value \\
    \midrule
    Vertex welding & round to \(10^{-6}\), replace by the group mean \\
    Small component & area below 1\% of the total surface \\
    Surface samples & 20{,}000, area-weighted, seed 0 \\
    Visibility directions & 64, Fibonacci sphere \\
    Ray offset & \(10^{-4}\) \\
    Thin-wall tolerances & 0.005 and 0.010 \\
    Sliver & triangle quality \(q < 0.1\) \\
    Dihedral exceedances & \(30^\circ\), \(60^\circ\), \(90^\circ\), length-weighted \\
    Decimation & quadric edge collapse to 10{,}000 faces; preserve topology, boundary and normals; planar quadrics \\
    Classifier & L2 logistic regression, \(C = 1\), balanced class weights, at most 5{,}000 iterations \\
    Threshold & MCC-optimal on 5-fold stratified out-of-fold scores, seed 0 \\
    Training assets & 548 silver-holdout meshes; asset 609 excluded \\
    \bottomrule
  \end{tabular}
\end{table}
% src: geomcheck/probes.py; analysis/PREREGISTRATION.md; results/RESULTS_milestone1.md

The primary decision thresholds locked by that rule, and then applied unchanged to every later corpus, are 0.416 for fused/incomplete, 0.440 for form/surface and 0.574 for extra geometry.
% src: results/ablations/lofo_family_rules.csv

Injection severities, in order S1, S2, S3, are: hole area fraction 0.002, 0.010, 0.040 across three patches; 1, 3 or 10 floaters of radius 0.01, offset by 0.05 along the normal; ellipsoid major radius 0.05, 0.10, 0.20 with transverse radii at half the major radius; fin thickness 0.008, 0.004, 0.002 on a plate of extent \(0.2 \times 0.2\); normal noise of standard deviation 0.001, 0.003, 0.010; duplicated patch area fraction 0.005, 0.020, 0.050 rotated by \(30^\circ\); hidden-shell radius 0.25, 0.50, 0.90 times the distance from the chosen interior point (the farthest from the surface among 20{,}000 uniform candidates) to the surface.
The per-object seed is the SHA-256 of the object name, reduced modulo \(2^{32}\).
The subset of 120 scans was drawn with \texttt{random.Random(0)} from the 1{,}033 Fuel models.
% src: scripts/gso_inject.py; analysis/PREREGISTRATION_M2.md

Feature families used in the ablation are those fixed in the milestone-2 plan: topology (component count, largest-component area fraction, small-component count and area, large-component count, genus proxy); boundary and manifold (boundary-edge fraction and length, boundary loops, winding inconsistency, duplicate-face fraction); self-intersection (face and area fractions); hidden surface; thin walls (defined-thickness fraction, both thin fractions, median thickness); roughness (dihedral mean and three exceedances, both angle-defect summaries); triangle quality (median, 5th percentile, both sliver fractions, edge-length coefficient of variation); size (face count).
% src: analysis/PREREGISTRATION_M2.md
